\newpage
\chapter{KESIMPULAN DAN SARAN} \label{Bab V}

\section{Kesimpulan} \label{V.Kesimpulan}
Berdasarkan hasil penelitian yang telah dilakukan, dapat disimpulkan beberapa hal sebagai berikut. \par

\begin{enumerate}[noitemsep]
	\item Pemanfaatan pendekatan spasiotemporal berbasis VideoMAE memberikan kontribusi positif dibandingkan pendekatan spasial berbasis Vision Transformer dalam klasifikasi video \textit{deepfake} pada dataset \textit{WildDeepfake}. Hal ini ditunjukkan oleh hasil eksperimen \textit{baseline}, yaitu model spasiotemporal BASE-ST memperoleh \textit{accuracy} sebesar 85,8\%, \textit{F1-score} sebesar 84,5\%, dan \textit{AUC} sebesar 92,3\%, sedangkan model spasial BASE-SP memperoleh \textit{accuracy} sebesar 79,2\%, \textit{F1-score} sebesar 76,7\%, dan \textit{AUC} sebesar 86,8\%. Peningkatan performa tersebut menunjukkan bahwa hubungan antar-\textit{frame} dalam satu klip memberikan informasi tambahan yang membantu model dalam membedakan klip \textit{real} dan \textit{fake} pada evaluasi internal \textit{WildDeepfake}. Dengan demikian, pada konteks dataset sumber, pendekatan spasiotemporal lebih efektif dibandingkan pendekatan spasial.
	\item Pendekatan spasial dan spasiotemporal menunjukkan sensitivitas yang berbeda terhadap perubahan karakteristik masukan dan konfigurasi pelatihan. Pada konfigurasi pelatihan, peningkatan \textit{learning rate} menjadi 5e-4 dan penggunaan augmentasi yang lebih berat cenderung menurunkan performa kedua pendekatan. \textit{Focal Loss} dan \textit{dropout} memberikan dampak yang lebih menguntungkan pada model spasial, tetapi tidak memberikan peningkatan pada model spasiotemporal. Sementara itu, \textit{CosineAnnealingLR} dipilih sebagai konfigurasi lanjutan karena memberikan rata-rata \textit{accuracy} yang lebih baik untuk kedua pendekatan. Pada konfigurasi input, perubahan resolusi, jumlah \textit{frame}, dan \textit{stride} juga memengaruhi performa model. Penurunan resolusi membatasi informasi visual yang dapat dipelajari model, sedangkan perubahan jumlah \textit{frame} dan \textit{stride} memengaruhi cakupan serta kepadatan informasi temporal yang digunakan dalam klasifikasi.
	\item Efektivitas pendekatan spasial dan spasiotemporal pada dataset \textit{WildDeepfake} tidak sepenuhnya konsisten ketika diterapkan pada dataset eksternal \textit{Celeb-DF v2} dalam skenario \textit{zero-shot}. Pada evaluasi lintas-\textit{dataset}, model spasial E6-SP memperoleh \textit{accuracy} sebesar 72,9\% dan \textit{AUC} sebesar 77,7\%, sedangkan model spasiotemporal E6-ST memperoleh \textit{accuracy} sebesar 64,2\% dan \textit{AUC} sebesar 72,6\%. Hasil ini menunjukkan bahwa pendekatan spasiotemporal yang lebih unggul pada evaluasi internal \textit{WildDeepfake} tidak selalu menghasilkan generalisasi terbaik pada dataset eksternal. Dengan demikian, fitur temporal yang efektif pada dataset sumber belum tentu memiliki transferabilitas yang sama pada distribusi data yang berbeda. Sebaliknya, pendekatan spasial berbasis Vision Transformer menunjukkan kemampuan generalisasi yang lebih baik pada skenario lintas-\textit{dataset}.
\end{enumerate}

Secara keseluruhan, penelitian ini menunjukkan bahwa perbandingan antara pendekatan spasial dan spasiotemporal perlu dilihat berdasarkan konteks evaluasinya. Pendekatan spasiotemporal berbasis \textit{VideoMAE} lebih efektif pada evaluasi internal \textit{WildDeepfake} karena mampu memanfaatkan hubungan antar-\textit{frame} secara langsung. Namun, pendekatan spasial berbasis Vision Transformer menunjukkan performa yang lebih stabil pada evaluasi lintas-\textit{dataset} terhadap \textit{Celeb-DF v2}. Oleh karena itu, pemilihan pendekatan untuk klasifikasi video \textit{deepfake} tidak hanya perlu mempertimbangkan performa pada dataset sumber, tetapi juga sensitivitas terhadap konfigurasi dan kemampuan generalisasi pada dataset eksternal. \par

\section{Saran} \label{V.Saran}
Berdasarkan hasil penelitian yang telah dilakukan, terdapat beberapa saran yang dapat dipertimbangkan untuk penelitian selanjutnya. \par

\begin{enumerate}[noitemsep]
	\item Penelitian selanjutnya disarankan untuk menerapkan \textit{fine-tuning} lintas-\textit{dataset}, atau \textit{multi-source training}, khususnya karena hasil evaluasi pada \textit{Celeb-DF v2} menunjukkan bahwa pendekatan yang unggul pada \textit{WildDeepfake} belum tentu memiliki transferabilitas yang sama pada \textit{dataset} eksternal. Langkah metodologis ini sejalan dengan rekomendasi Nadimpalli dan Rattani yang menegaskan bahwa adaptasi domain diperlukan untuk meningkatkan generalisasi detektor \textit{deepfake} ketika dihadapkan pada perbedaan distribusi data dunia nyata \cite{nadimpalli2022improvingcrossdatasetgeneralizationdeepfake}.
	\item Penelitian selanjutnya disarankan untuk mengeksplorasi strategi \textit{sampling} temporal yang lebih adaptif, seperti \textit{multi-clip sampling}, \textit{segment-based sampling}, atau \textit{saliency-based sampling}. Saran ini didasarkan pada hasil eksperimen konfigurasi input yang menunjukkan bahwa jumlah \textit{frame} dan \textit{stride} memengaruhi performa model. Strategi pemilihan klip yang dinamis ini didukung oleh temuan Korbar et al. melalui \textit{SCSampler}, yang menunjukkan bahwa pemilihan klip video yang lebih salien dapat mengurangi redundansi pemrosesan dan meningkatkan efisiensi pengenalan video \cite{korbar2019scsamplersamplingsalientclips}.
	\item Penelitian selanjutnya juga dapat mengembangkan pendekatan gabungan antara model spasial dan spasiotemporal, misalnya melalui \textit{late fusion} pada tingkat \textit{logit} atau probabilitas, maupun integrasi fitur sebelum tahap klasifikasi. Pendekatan ini berpotensi memanfaatkan keunggulan model spasial dalam menangkap artefak visual \textit{frame-level} dan keunggulan model spasiotemporal dalam memodelkan hubungan antar-\textit{frame}. Konsep integrasi fitur ruang dan waktu ini diperkuat oleh penelitian Kingra et al. yang merancang arsitektur \textit{SFormer}, yaitu kombinasi antara analisis spasial menggunakan Transformer dan analisis temporal menggunakan blok lanjutan. Pendekatan tersebut terbukti menghasilkan deteksi video \textit{deepfake} yang lebih komprehensif serta lebih \textit{robust} terhadap pergeseran distribusi data eksternal \cite{KINGRA2024301817}.
\end{enumerate}
